\documentclass[11pt]{article}
\usepackage[margin=1in]{geometry}
\usepackage{listings}
\usepackage{xcolor}
\usepackage{graphicx}
\usepackage{hyperref}

\definecolor{codegray}{rgb}{0.5,0.5,0.5}
\definecolor{codegreen}{rgb}{0,0.6,0}
\definecolor{backcolour}{rgb}{0.95,0.95,0.92}

\lstdefinestyle{mystyle}{
  backgroundcolor=\color{backcolour},
  commentstyle=\color{codegreen},
  keywordstyle=\color{magenta},
  numberstyle=\tiny\color{codegray},
  basicstyle=\ttfamily\footnotesize,
  breaklines=true,
  numbers=left,
  numbersep=5pt,
  showstringspaces=false,
  tabsize=2,
  frame=single
}
\lstset{style=mystyle}

\title{CS 2110 Project 2: Building an LRU Cache}
\author{Ethan Brooks \and Aisha Rahman}
\date{\today}

\begin{document}
\maketitle

\section{Overview}
A least-recently-used (LRU) cache evicts the entry that has gone unused for the longest
time once its capacity is reached. Both \texttt{get} and \texttt{put} must run in $O(1)$
time. We combine a hash map with a doubly linked list.

\section{Implementation}
Listing~\ref{lst:lru} shows our Python implementation.

\begin{lstlisting}[language=Python, caption={LRU cache with an ordered dictionary.}, label={lst:lru}]
from collections import OrderedDict

class LRUCache:
    def __init__(self, capacity: int):
        self.capacity = capacity
        self.data = OrderedDict()

    def get(self, key):
        if key not in self.data:
            return -1
        self.data.move_to_end(key)   # mark as recently used
        return self.data[key]

    def put(self, key, value):
        if key in self.data:
            self.data.move_to_end(key)
        self.data[key] = value
        if len(self.data) > self.capacity:
            self.data.popitem(last=False)  # evict oldest
\end{lstlisting}

\section{Testing}
We tested the cache with the sequence below; inline code such as
\lstinline{cache.get(1)} returns \lstinline{-1} after eviction.

\begin{lstlisting}[language=Java, numbers=none]
LRUCache cache = new LRUCache(2);
cache.put(1, 1);
cache.put(2, 2);
cache.get(1);       // returns 1
cache.put(3, 3);    // evicts key 2
\end{lstlisting}

\section{Complexity}
\begin{tabular}{|l|c|c|}
\hline
Operation & Time & Space \\
\hline
\texttt{get} & $O(1)$ & --- \\
\texttt{put} & $O(1)$ & $O(n)$ total \\
\hline
\end{tabular}

\section{Conclusion}
The ordered dictionary gives a compact solution. See the course page at
\url{https://cs2110.example.edu/projects} for the grading rubric.

\end{document}
